% page 506 of the facsimile (image p-505 of the scan)
% block 160.0 x 233.3 mm ; left margin 8.0 mm
\pgc{-12.700mm}{0.000mm}{
\l{\cel{3}498.}
\l{}
\l{\sou{rezistar}. (\sou{trans.)}\cel{2}I. Opozar, a la efiko da forco, forcon}
\l{\cel{3}qua agas inversa-sinse. - II. Aplikar esforco kontre la}
\l{\cel{3}uzo di forco, di violento. - III. Opozar sua volo ad im-}
\l{\cel{3}pulso, a volo adversa. - EFIS.}
\l{}
\l{\sou{rezolvar}. (\sou{trans}.) Determinar, pos decido, to quo esas agenda}
\l{\cel{3}praktikale, aplikale. - DEFIS.}
\l{}
\l{\sou{rezonar}.\cel{2}(\sou{netrans.,}\cel{1}pri) Agar kateniguro ek propozicioni qui}
\l{\cel{3}esas apta demonstrar ulo. - DEFIRS.}
\l{}
\l{\sou{rezultar}.\cel{2}(\sou{netrans.,}\cel{1}de).Eventar konseque de ago, de fakto.}
\l{\cel{3}DEFIRS.}
\l{}
\l{\sou{rezumar}. (\sou{trans.}) Repetar to quo dicesis skribe o voce, ma}
\l{\cel{3}agar lo plu kurte, per expresar nur lo esencala o quitesen-}
\l{\cel{3}cala. - DEFIS.}
\l{}
\l{\sou{rezurektar}. (\sou{netrans.}) (\sou{religio kristana)}.\cel{2}Retrovenar de mort-}
\l{\cel{3}inteso a vivo, per miraklo. -}
\l{}
\l{\sou{ri-}\cel{1}.\cel{3}Prefixo qua indikas ago quan onu iteras, iteris od}
\l{\cel{3}iteros.}
\l{}
\l{\sou{ribo}. (\sou{bot.}) Frukto di arbusto ek la familio \textquotedbl{}grosulariei\textquotedbl{}.}
\l{\cel{3}- L. ribes rubrum. IL.}
\l{}
\l{\sou{richa}. Qua havas propriete multa havaji importoza. (richa}
\l{\cel{3}de). - DEFIS.}
\l{}
\l{\sou{ricino}. (bot.) Genero de planti, ek la familio \textquotedbl{}euforbiacei\textquotedbl{},}
\l{\cel{3}di qua un speco donas oleo purgiva. - L.\cel{1}\sou{ricinus}. DEFISL.}
\l{}
\l{\sou{\textquotedbl{}rictus\textquotedbl{}}. Fenduro di la boko, qua lasas la denti videsar.}
\l{}
\l{\sou{ridar}. (\sou{netrans., pri, pro,}\cel{1}de).Sentar quaza expanseso di sua}
\l{\cel{3}vizajo quan akompanas expiri sukusoza plu o min bruisanta,}
\l{\cel{3}signi korpala di gayesko subita. - FIS.}
\l{}
\l{\sou{ridelo}. (\sou{tekn.}) Ligno-peco horizontala qua, ye singla latero}
\l{\cel{3}di chareto, formacas, per la ronda ligno-cilindri qui tra-}
\l{\cel{3}iras lu, quaza rastelero destinita impedor la falo di la}
\l{\cel{3}kargajo. - FI.}
\l{}
\l{\sou{rifo.}\cel{2}(\sou{navig.})\cel{2}En maro, rokajo, grupo de rokaji o de sablo}
\l{\cel{3}o de koralii, nur kelke salianta de la aquo-surfaco, kontre}
\l{\cel{3}qua povas su ruptar navi. - DER.}
\l{}
\l{\sou{rigo}. I. (\sou{navig.)}\cel{3}To quo uzesas por manovrar navo (kordegi,}
\l{\cel{3}segli, e c.) od aerostato. - II. Singla de la kordegi,}
\l{\cel{3}korda skali, ligna stangi qui pendas de la portiko digim-}
\l{\cel{3}nastikeyo. - E.}
}
